\section{Discussion}
\label{sec:discussion}

\subsection{Computational efficiency}
A full 600-iteration optimisation on the primary $80\times20$ mesh
completes in under thirty minutes on a standard laptop (Intel Core i7,
16\,GB RAM).
Each iteration requires five sparse linear systems (two forward, two
adjoint, one Helmholtz filter), all handled by MUMPS direct factorisation.
For larger meshes the iterative block-preconditioner (Section~\ref{sec:fem})
maintains near-optimal convergence rates (Supplementary~S1).

\subsection{Mesh resolution and optimizer sensitivity}
\label{sec:mesh_limitation}
On coarse meshes ($20\times5$, $N_{\mathrm{out}}=6$) the outlet
discretisation is insufficient to represent the linear gradient target
and the optimizer finds unphysical reversed-flow local minima.
The $40\times10$ mesh ($N_{\mathrm{out}}=11$) converges to a
suboptimal basin; systematic improvement requires the $80\times20$
resolution or finer.
On the primary mesh, the transverse diffusion layer thickness
$\delta \approx 4\,\mu$m (Eq.~\eqref{eq:delta}) is underresolved by
a factor of $\approx 6\times$ ($\Delta y = 25\,\mu$m).
SUPG eliminates spurious oscillations but does \emph{not} recover
sub-grid diffusion structure.
The RMSE metric measures the \emph{global} outlet profile over the full
channel width $L_y = 500\,\mu$m, a scale that is well resolved;
the channel-scale branching topology is therefore correctly captured
even though the wall-normal diffusion sublayer is not.

Mesh convergence (Table~\ref{tab:mesh_conv}, Supplementary~S1) confirms
monotone RMSE decrease with refinement.
Full convergence to within 1\% of the continuum value requires
$\Delta y \lesssim \delta \approx 4\,\mu$m, demanding adaptive mesh
refinement or a mesh of $\gtrsim 500\times125$ elements;
this is deferred to future work.
With an identical 600-iteration hybrid OC/MMA schedule, the $160\times40$ mesh gives RMSE~$=\rmseHiRes=0.091$ and $J=9.82\times10^{-4}$, compared to RMSE~$=\rmseTopOpt=0.064$ and $J=1.77\times10^{-3}$ on the $80\times20$ mesh.
The RMSE is non-monotone (finer mesh gives higher outlet RMSE), while the objective $J$ IS monotone ($J_{160}=9.82\times10^{-4} < J_{80}=1.77\times10^{-3}$), suggesting the finer mesh finds a lower-dissipation topology with higher gray fraction (0.676 vs 0.540) that achieves better flow efficiency but less sharp concentration gradients.
True Richardson-extrapolation mesh convergence would require identical topology similarity ($\|\rho_h - I(\rho_{2h})\|/\|\rho_{2h}\|\to 0$) and is deferred to future work; the current study demonstrates optimizer sensitivity to mesh resolution rather than pure discretisation error.
          \emph{The Brinkman surrogate RMSE~$=\rmseTopOpt$ is
          not the performance of a physically realisable chip};
          closing this gap remains an open challenge.
    \item \textbf{Brinkman approximation.}
          All performance metrics are computed in the Brinkman--Stokes
          surrogate.
          At $\Da \approx 4\times10^{-6} \ll 1$ in solid regions, the
          surrogate accurately represents impermeability~\cite{borrvall2003}.
          Body-fitted Navier--Stokes validation is implemented in
          \texttt{ns\_validation.py} but requires \texttt{pyvista} and
          \texttt{gmsh} dependencies not present in the standard Docker
          image; these results are deferred to a companion study.
    \item \textbf{Two-dimensional steady flow.}
          The framework is 2D; three-dimensional extension is implemented
          in \texttt{validation\_3d.py} but not systematically validated.
    \item \textbf{Fixed objective weights.}
          The weights $w_f$ and $w_c$ are fixed in this study.
          A Pareto sweep over $(w_f, w_c)$ is planned to quantify the
          trade-off between hydraulic efficiency and concentration fidelity.
    \item \textbf{No experimental validation.}
          Results are purely computational.
          Experimental fabrication and testing are necessary next steps.
\end{itemize}

\subsection{Future directions}
Immediate priorities are: (i)~adaptive mesh refinement targeting the
fluid--solid interface; (ii)~integration of body-fitted Navier--Stokes
post-processing into the main pipeline; (iii)~a systematic Pareto front
study over objective weights; (iv)~three-dimensional validation;
and (v)~experimental collaboration for a representative subset of designs.
